% Citation pass: the SSIM and Hounsfield-unit entries were re-scoped from what the
% tomographic chapter reports to what it frames and states, so that after the arXiv-only
% reconstruction records were withdrawn no entry describes a measured operating point, or a
% reported error value, that the book no longer carries.
% GLOSSARY -- the chapter title and the label app:glossary are fixed, since Ch. 3.8
% cross-references app:glossary and Ch. 1.1/1.5/2.1/2.5/2.6
% reference the appendices this chapter sits beside. Every entry is a definition
% of the term AS THIS BOOK USES IT, drawn from three places and no others:
% the master notation table (frontmatter/notation.tex, tab:notation-master), the
% definition and assumption environments of Ch. 1.1--1.9 (def:generic_no,
% def:ch06-discretization-invariance, ass:ellipticity, ass:ch09-compact), and the
% working usage of Parts II and III. No entry states a result, a number, or a
% bound that is not printed in the chapters it cites; where a chapter reports a
% quantity without defining it (SSIM, PSNR) the entry says so rather than
% supplying an outside formula. Index anchors are added only where they do not
% duplicate a bold anchor already set in the main text, and reuse the main text's
% exact key strings so that makeindex merges rather than splits them.
\chapter{Glossary}\label{app:glossary}

This glossary collects the vocabulary the monograph actually argues with: the
terms of Part~\ref{part:1}, the tensor and runtime names of Part~\ref{part:2}, and
the physical and clinical quantities of Part~\ref{part:3}. It is a consolidation,
not a second text. Each entry states what the book means by its term, gives the
display, definition, or table in which that meaning is fixed, and records the
qualification the chapter attaching the term insists on. Where the chapters use a
word without defining it, the entry says so; where two chapters use one letter for
two objects, the entry names both and points to Table~\ref{tab:notation-master},
which governs. Nothing here is a pointer: a reader who has met neither neural
operators nor the discretizations they consume should be able to read an entry
and know what the term denotes, what it is bounded by, and what it is not.

Two naming conventions are worth stating before the entries, because the book is
deliberate about both. First, the architecture names are the ones the papers of
record use, and each is expanded on first use in the part where it appears;
Appendix~\ref{app:e} gives the variant-by-variant inventory and
the chapter in which each is introduced, and this glossary does not duplicate it.
Second, error vocabulary is not interchangeable across the three parts. The
triangle inequality of \eqref{eq:ch01-error-split} names approximation,
discretization, and statistical error; \eqref{eq:ch04-error-split} names sensor,
basis, and network error; \eqref{eq:ch05-error-decomposition} names spectral
truncation and neural parametric error; \eqref{eq:ch06-decomp} names statistical,
truncation, and discretization error in the three-way form
of~\cite{subedi2024controlling}. Table~\ref{tab:appg-error} sets those side by
side, and the conclusion to take from it is the one \chapref{ch:part1-introduction-the-paradigm-shift-to-infinit}
draws: the terms are charged to different variables and decay at different rates,
so a single ``relative error'' measured at the training resolution confounds the
approximation term with the discretization term and may describe the quadrature
rather than the learned operator.

\begin{table}[t]
\centering
\small
\caption{\textbf{The book's error vocabulary, consolidated.} Each row is a term as
displayed in the cited equation; ``charged to'' names the design variable the chapter
attaches it to, and the last column reports only what that chapter states. The seven
rows are not seven terms: different parts cut the same total error along different
lines, and the final column records where each cut is priced.}
\label{tab:appg-error}
\begin{tabularx}{\textwidth}{@{}l >{\raggedright\arraybackslash}X l >{\raggedright\arraybackslash}X@{}}
\toprule
\textbf{Term} & \textbf{Displayed in} & \textbf{Charged to} & \textbf{What the book states} \\
\midrule
Statistical & \eqref{eq:ch01-error-split}, \eqref{eq:ch06-decomp}, \eqref{eq:ch09-rate} & $n$ & $O(n^{-2s/(2s+d)})$ for the risk \eqref{eq:ch09-risk} under Assumption~\ref{ass:ch09-compact}; no bound of that form exists when $s = 0$, and a sharp form for the operator classes this book learns is listed as open \\
Approximation & \eqref{eq:ch01-error-split}, \eqref{eq:ch05-error-decomposition} & $L$, $d_{v}$ & Theorem~\ref{thm:no_universality} asserts existence and is silent on how large depth, width, and the kernel parameters must grow as $\varepsilon \to 0$ \\
Discretization & \eqref{eq:ch01-error-split}, \eqref{eq:ch06-rate} & $N$ & bounded by $C N^{-s}$ at fixed $\theta$ for inputs in $H^{s}$, $s > d/2$; the mechanism is aliasing \\
Truncation & \eqref{eq:ch05-error-decomposition}, \eqref{eq:ch06-decomp} & $k_{\max}$ & $C_{1} k_{\max}^{-\alpha}$, with $\alpha$ the smoothing order of the target; set by regularity, not by depth, width, or grid \\
Sensor & \eqref{eq:ch04-error-split}, \eqref{eq:ch8-sensor-rate} & $m$ & $\lesssim m^{-s/d}$ for ensembles bounded in $H^{s}$, $s > d/2$; recorded as the expectation covering asymptotics justify, not as a bound for a trained network \\
Basis & \eqref{eq:ch04-error-split} & $p$ & projection error onto the $p$-dimensional span of the trunk outputs; shrinks quickly for smooth outputs, no rate stated \\
Network & \eqref{eq:ch04-error-split}, and the term \eqref{eq:ch01-error-split} omits & $\theta$ & the only term a longer training run reduces without enlarging $m$ or $p$; optimization error is treated empirically in Part~\ref{part:2} \\
\bottomrule
\end{tabularx}
\end{table}

The entries follow. They are alphabetized by the headword as printed, and each
carries the index anchor under which the reader will find it in the subject index.

\begin{description}

\item[\textbf{Aliasing}] The contamination of a resolved Fourier coefficient by
every coefficient whose mode differs from it by a multiple of the grid size $N$,
which is exactly what the discrete transform of nodal values computes in place of
the continuous integral; the identity is \eqref{eq:ch06-alias}, and
\chapref{ch:part1-discretization-invariance-and-continuum-li} names aliasing as the
mechanism of the discretization bound \eqref{eq:ch06-rate}. The same arithmetic
reappears on the input side: when a sensor ensemble carries energy above the
sensor Nyquist rate, the sampled representation folds it back onto the resolved
coefficients, so the branch is trained on a corrupted encoding rather than an
incomplete one, and neither a wider trunk nor a longer run repairs it.
\index{aliasing}

\item[\textbf{Anomaly correlation coefficient (ACC)}] The verification score of
\eqref{eq:ch33-acc}: the correlation between forecast and analysis \emph{anomalies},
each taken as the departure from a climatological field $\mathbf{C}$. Because
climatology is subtracted from both sides, the coefficient is blind to additive
bias, and a forecast displaced uniformly from the analysis still scores one, which
is why Part~\ref{part:3} always quotes it beside the root-mean-square error. As
written it is a sum over array indices, so on a latitude--longitude sphere it
weights a polar square kilometer more heavily than an equatorial one; the remedy is
the area-weighted form $\mathrm{ACC}_{w}$ of \eqref{eq:ch33-accw} with $w_j =
\cos\phi_j$, which the published verification uses. A forecast horizon is whatever
threshold of one agrees to cross; the book reports ACC above $0.6$ for total column
water vapor beyond day eight as that verification's own claim.
\index{anomaly correlation coefficient (ACC)|textbf}

\item[\textbf{Automatic mixed precision (AMP)}] The PyTorch facility, entered
through \texttt{torch.cuda.amp.\allowbreak autocast}, that runs part of a network in a
16-bit format; \chapref{ch:part2-training-regimes-optimization-and-failure-} evaluates it
against the formats of Table~\ref{tab:ch26-formats} and reaches a split verdict.
\texttt{bfloat16} has eight exponent bits, so it overflows or underflows only where
\texttt{float32}$'$s dynamic range already fails, but its unit roundoff $2^{-8}$ is
charged once per butterfly stage, and the per-stage charge is levied inside the
transform on quantities every later stage reads; \texttt{float16} fails additionally
on range and needs loss scaling in the backward pass. The policy the chapter states
is to run the pointwise $1 \times 1$ lifting and projection convolutions in
\texttt{bfloat16} and keep the transform and the complex mode contraction in
\texttt{float32}, expressed as an explicit cast at the boundary of the spectral
layer because \texttt{autocast} assigns dtypes by operation type and no operation
table records which sum is the measured one.
\index{automatic mixed precision (AMP)|textbf}

\item[\textbf{Autoregressive rollout}] The recursion $v_{k+1} =
\mathcal{G}_{\theta}(v_k)$ of \eqref{eq:ch26-rollout}, in which the trained object is
a time-$\Delta t$ map applied to its own output, with $v_0 = u_0$. Part~\ref{part:1}
prices it in \eqref{eq:ch09-onesided} and \eqref{eq:ch09-gronwall} through the
one-step error $\tau$ and the network's global Lipschitz constant $\Lambda_\theta$,
and the arithmetic is unforgiving: if $\Lambda_\theta$ exceeds one, the bound grows
geometrically in the number of steps however small $\tau$ is. The book's regimes are
not corner cases --- the reference configuration unrolls the two-dimensional
Navier--Stokes vorticity equation across fifty to one hundred steps, and the weather
surrogate of \cite{pathak2022fourcastnet} advances six-hour steps over days. Because
a rollout is queried at the model's own iterates rather than at solver snapshots,
the distribution of training inputs and of inference inputs diverge, which is the
defect \eqref{eq:ch26-pushforward} addresses. \index{autoregressive rollout|textbf}

\item[\textbf{Banach space}] A normed vector space complete for the metric its norm
induces, so that every Cauchy sequence converges to an element of the space;
completeness is the hypothesis that makes the limit of an approximating sequence
again a field with its norm bound preserved. For the spaces of this book it is a
theorem rather than an assumption: the Riesz--Fischer theorem gives $L^{p}(\Omega)$
complete for every $1 \le p \le \infty$, and Sobolev spaces are closed subspaces of
finite products of them. Where the norm comes from an inner product the space is
Hilbert, and the distinction is substantive --- among the Lebesgue spaces only
$p = 2$ satisfies the parallelogram law, so orthogonality, best approximation from a
subspace, and Parseval's identity are available in $L^{2}$ and unavailable in
$L^{1}$ or $L^{\infty}$. The compactness failure that organizes Part~\ref{part:1} is
also a Banach-space fact: by Riesz's lemma the closed unit ball is compact exactly
when the space is finite-dimensional. \index{Banach space|textbf}

\item[\textbf{Branch network}] The DeepONet sub-network that maps the $m$ sensor
readings of the input function to a latent code $\mathbf{b}(u) \in \mathbb{R}^{p}$,
the first factor of the contraction \eqref{eq:ch04-deeponet} and the term that
carries the encoding error $\mathcal{E}_{\text{sensor}}(m)$ of
\eqref{eq:ch04-error-split}. Its first layer consumes $m$ numbers in a fixed order,
so the sensor locations enter the architecture only through their count --- one
branch pass serves every query point, which is why the input side is the sensor-bound
side. In the PyTorch module of \chapref{ch:part2-deeponet-engineering-branch-trunk-networks}
its output width is the shared latent width \texttt{latent\_p}, and the constructor
welds $m$ and $p$ into the parameter tensors before the first epoch, so neither is a
hyperparameter that a later sweep can revise. \index{DeepONet!branch network|textbf}

\item[\textbf{Burgers' equation}] $\partial_{t} u + u\,\partial_{x} u =
\nu\,\partial_{xx} u$ on $x \in [0,1)$, $t \in [0,1]$ with periodic boundary
conditions \eqref{eq:ch25-burgers}, so the domain is the torus
$\mathbb{T} = \mathbb{R}/\mathbb{Z}$ and not an interval --- periodicity both
licenses the Fourier representation and removes any boundary condition from the
specification, which makes this the cleanest of the book's three tests of
discretization invariance. The state is scalar, so input and output channel counts
are both one, and the target is the forward time-evolution operator
$u_{0} \mapsto u(\cdot, t)$ of \eqref{eq:ch03-burgers}, whose domain and codomain
coincide. At $\nu > 0$ the map smooths, sending $L^{2}$ data into $H^{s}$ for every
$s$; the figure at $\nu = 0.01$ shows a settled viscous shock well past the inviscid
gradient catastrophe at $t_{*} \approx 0.16$. The maximum principle
$\sup_{x} u(\cdot, t) \le \sup_{x} u_{0}$ is used in Part~\ref{part:2} as an
implementation check: a computed profile exceeding $\|u_{0}\|_{L^{\infty}}$ has
produced an artifact, and no relative $L^{2}$ number will tell you so.
\index{Burgers equation|textbf}

\item[\textbf{Compact embedding; compact operator}] A linear operator between normed
spaces is compact when the image of every bounded set has compact closure,
equivalently when bounded sets are sent to precompact ones; the canonical injection
$\iota : H^{s}(\Omega) \hookrightarrow H^{t}(\Omega)$ of
\eqref{eq:ch02-injection} is the book's named instance, asserted for $s > t \ge 0$
on bounded $\Omega$ with $C^{1}$ boundary by Theorem~\ref{thm:ch02-rellich}. The
notion is load-bearing rather than decorative: because no closed bounded set in an
infinite-dimensional space is compact, the uniform approximation theorems of
Part~\ref{part:1} quantify over compact sets, and Rellich--Kondrachov is the
machinery that converts a regularity bound into one. On the torus the content is the
tail estimate \eqref{eq:ch02-tail}, whose rate $(1+K^{2})^{(t-\sigma)/2}$ is the
algebraic truncation decay the Fourier neural operator's mode cutoff inherits.
\index{compact embedding|textbf}

\item[\textbf{Continuously differentiable spaces $C^{k}$}] $C^{k}(\overline{\Omega})$,
the $k$-times continuously differentiable functions with the uniform norm, is the
target the notation table lists alongside $L^{p}$ and $H^{s}$, and the space in which
pointwise statements are made. Two results carry the book between the integral and
the pointwise: the Arzel\`a--Ascoli theorem characterizes relatively compact subsets
of $C(\overline{\Omega})$ as the uniformly bounded equicontinuous ones, and the
supercritical clause of Proposition~\ref{prop:ch02-sobolev-embeddings} embeds
$W^{k,p}$ into $C^{j,\alpha}$ whenever $j + \alpha \le k - d/p$, the only route by
which an integral bound implies a pointwise one. \chapref{ch:part1-operator-approximation-theory-and-universa}
uses the second to place the Darcy solution in $C(\overline{\Omega})$ once the
available Sobolev order exceeds $d/2$.
\index{continuously differentiable spaces|textbf}

\item[\textbf{Coordinate channels (positional grid encoding)}] The normalized spatial
coordinates appended as extra input channels, $a_{\text{augmented}}(x) =
[a(x), x_{1}, \dots, x_{d}] \in \mathbb{R}^{C_{\text{in}} + d}$ of
\eqref{eq:ch21-augment}, concatenated along the channel axis and therefore requiring
the pointwise lift $\mathcal{P}$ to be resized for $C_{\text{in}} + d$ inputs. The
extra channels are not data but the argument of the function, and they are needed
because the spectral branch commutes with translation --- a shift multiplies each
coefficient by a phase the multiplier never compares with anything --- so a network
built from stationary kernels and pointwise maps cannot express a preference for
absolute location. The price is a loss of regularity: the coordinate ramp is not
periodic on the torus the transform assumes, so augmented inputs lie in
$H^{1/2 - \epsilon}$ for every $\epsilon > 0$ and no better, after which the
discretization rate \eqref{eq:ch06-rate} never reaches $N^{-1}$. The density results
neither require nor forbid them; mesh layouts, where the coordinates already exist as
a separate tensor and the kernels consume them as arguments, need no such
augmentation. \index{coordinate channels|textbf}\index{positional grid encoding}

\item[\textbf{Courant--Friedrichs--Lewy (CFL) condition}] The stability restriction on
an explicit time-marching scheme, $\Delta t \le \mathrm{CFL}\,\Delta x /
|\mathbf{u}|_{\max}$, which \eqref{eq:ch31-refine} uses to show that refining a grid
raises the unknown count $N$ and the step count $K$ together, so one classical run
costs $NK = \mathcal{O}(N^{1+1/d})$. That coupling is a premise of the multi-query
argument of \chapref{ch:part3-the-multi-query-bottleneck-in-computationa}, and it
reappears as a cost of data generation. The reference solvers of
Part~\ref{part:2} take their own step small enough to satisfy the constraint, and the
cardiac counterpart \eqref{eq:ch35-cfl} is the diffusion restriction
$\Delta t \le h^{2}/(2 d D_{\max})$, where $D_{\max}$ bounds the eigenvalues of the scaled
diffusion tensor $\mathcal{D} = \mathbf{D}/(\chi C_{m})$ rather than any advective speed.
\index{Courant--Friedrichs--Lewy condition|textbf}

\item[\textbf{Darcy flow}] Single-phase saturated porous flow: mass conservation
$\nabla \cdot \mathbf{q} = f$ with Darcy's empirical flux law
$\mathbf{q} = -a\nabla u$ gives the steady diffusion problem
$-\nabla \cdot (a(x)\nabla u(x)) = f$ with $u|_{\partial\Omega} = 0$
\eqref{eq:darcy}, where $a \in L^{\infty}(\Omega)$ is the permeability coefficient
and $u$ the pressure potential. It is the book's model \emph{coefficient-to-state}
map: input and output spaces differ, the energy estimate lands the solution in
$H^{1}_{0}(\Omega)$, and smooth coefficients lift it to $H^{2}(\Omega)$, so the
operator is smoothing rather than evolutionary. The benchmark
\eqref{eq:ch25-darcy} writes the coefficient $K$, poses the problem on the unit
square with homogeneous Dirichlet data, and thereby forfeits the periodicity a
spectral architecture is exact on --- a specification that does not say whether the
mismatch is absorbed or the field odd-extended leaves its dominant error term
unnamed. Assumption~\ref{ass:ellipticity}, uniform ellipticity
$0 < a_{\min} \le a \le a_{\max} < \infty$, is what makes the elliptic estimate
uniform, and \chapref{ch:part3-subsurface-multiphase-flow-and-carbon-stor} generalizes
the law to two phases \eqref{eq:ch34-flux}, where each phase carries a relative
permeability $k_{r\alpha}(S_{\alpha}) \in [0,1]$ and the phases are coupled by
saturation completeness and a capillary pressure. \index{Darcy flow|textbf}

\item[\textbf{DeepONet (Deep Operator Network)}] The architecture
$\mathcal{G}_{\theta}(u)(y) = \sum_{k=1}^{p} b_{k}(u(x_{1}), \dots,
\allowbreak u(x_{m}))\,\allowbreak t_{k}(y) + b_{0} = \langle \mathbf{b}(u),
\allowbreak \mathbf{t}(y)\rangle_{\mathbb{R}^{p}}\allowbreak + b_{0}$ of \eqref{eq:ch04-deeponet},
which turns the separable rank-$p$ form of
Theorem~\ref{thm:chen_chen} into two trainable networks joined by one inner
product~\cite{lu2021learning}. Its defining feature is an asymmetry: the output side
is coordinate-free, since $y$ is an argument of the trunk and may be queried anywhere,
while the input side is tied to $m$ prescribed sensors. Theorem~\ref{thm:deeponet_error}
prices the two sides and the network between them separately
\cite{lanthaler2022error}, and \eqref{eq:ch8-sensor-rate} gives the encoding term the
decay $m^{-s/d}$ that a bounded input ensemble in $H^{s}$, $s > d/2$, justifies.
\index{DeepONet}

\item[\textbf{Discretization invariance}] The property of an \emph{architecture
family}, not of a trained model, stated in Definition~\ref{def:ch06-discretization-invariance}
as three conditions: the parameter space is defined independently of any mesh, each
parameter determines a continuum operator between function spaces, and every finite
discretization induces a discrete operator whose diagram with the restriction
$P_N$ and extension $I_N$ commutes in the continuum limit, as
\eqref{eq:limit_commute} and Figure~\ref{fig:commutative} express. The conditions
fail separately --- a densely connected network on a fixed grid violates the first
trivially --- and the third hides the analysis, since it requires point evaluation to
be bounded on the input space, which on $L^{2}(\Omega)$ it is not. Two omissions are
part of the meaning: \eqref{eq:limit_commute} is a limit statement with $\theta$
fixed, asserting existence of a continuum operator and not a rate, and it says
nothing whatsoever about the parameters a training procedure returns when it is run
on a coarse grid. \index{discretization invariance|textbf}

\item[\textbf{Duality pairing}] The evaluation of a continuous linear functional on a
Banach space, written $\langle \ell, v \rangle_{X^{*}, X}$ for $\ell \in X^{*}$ and
$v \in X$; the notation keeps the two spaces visibly distinct, unlike an inner product,
which is what lets a source term that is a distribution rather than a function still be
tested against smooth functions. The pairing is defined through the operator norm
\eqref{eq:ch02-dualnorm}, a supremum over the unit ball, so it measures worst-case
amplification. The book's use of it is a choice of norm, and the choice changes what
counts as small: replacing $L^{2}$ by $H^{-1}$ makes the residual of a differential
equation far easier to make small, \eqref{eq:ch02-hsneg} extends Parseval's identity to
that weaker scale, and the pairing extends the inner product of the pivot space
\emph{downward} in regularity through the triple $V \hookrightarrow H \cong H^{*}
\hookrightarrow V^{*}$ of \eqref{eq:ch02-triple}. A residual small in $H^{-s}$ is small
against every test field of unit $H^{s}$ mass, a strictly weaker demand than smallness in
$L^{2}$. \index{duality pairing|textbf}

\item[\textbf{Einstein summation convention (\texttt{einsum})}] The labeled-index
notation for tensor contractions, and the PyTorch operation \texttt{torch.einsum}
that executes it. In this book it names one specific contraction at a time: the
per-mode channel mixing \texttt{"bixy,ioxy->boxy"} of
\chapref{ch:part2-fourier-neural-operators-fno-from-scratch}, applied to the two
contiguous corners of the retained half-spectrum against the complex weight tensors
$R^{(1)}, R^{(2)}$ of \eqref{eq:ch23-Rtensors}. Two cautions come with it. The
contraction's cost carries no factor of the node count, so refinement dilutes the share
of the cost that any low-rank factorization of the multiplier array could take; and a letter
repeated inside one operand takes the diagonal of the two axes it labels, so it raises a shape error
when their lengths differ and returns wrong numbers without complaint when they
coincide --- the printed DeepONet contraction
of \chapref{ch:part2-deeponet-engineering-branch-trunk-networks} uses \texttt{p} for
both the latent and the query axis and performs the intended operation of
\eqref{eq:ch22-batch-contraction} only when the two are labeled apart.
\index{Einstein summation convention (einsum)|textbf}

\item[\textbf{ERA5 reanalysis}] The European Centre for Medium-Range Weather
Forecasts atmospheric reanalysis \cite{hersbach2020era5}, the source of the labels the
weather surrogate of \chapref{ch:part3-global-weather-forecasting-and-climate-dyn}
is fitted to: hourly snapshots subsampled to $00$, $06$, $12$, $18$\,UTC on a
$0.25^{\circ} \approx 28$\,km grid, holding the twenty prognostic variables of
\eqref{eq:ch33-statemap}. Two bookkeeping details matter for reproduction. The
archived regular grid carries $721$ latitude nodes because both poles are included,
so the tensor the code reads is $721 \times 1440 \times 20$; and ERA5's native
horizontal grid is a Gaussian grid, re-interpolated to the regular latitude--longitude
arrangement by the archive's distribution service before training. The book records
that the interpolation is a low-pass operation applied to the labels --- arithmetically
benign, physically not --- so any claim about what such a model resolves must be read
against it. \index{ERA5 reanalysis|textbf}

\item[\textbf{Error decomposition}] The practice of splitting a surrogate's total
error into terms that respond to different design variables. The book displays four
cuts: the three-term triangle inequality \eqref{eq:ch01-error-split}, the DeepONet
split \eqref{eq:ch04-error-split}, the spectral split
\eqref{eq:ch05-error-decomposition}, and the statistical, truncation, and
discretization form \eqref{eq:ch06-decomp}, made explicit with matching upper and
lower bounds for a discrete-Fourier-transform least-squares estimator of the linear
layer by \cite{subedi2024controlling}. The book records one term the first of these
displays omits: optimization error, the gap between the parameter stochastic gradient
descent returns and an exact minimizer. Part~\ref{part:2} treats it empirically, while the DeepONet split carries the same gap explicitly as its
network term $\mathcal{E}_{\text{DNN}}(\theta)$. Table~\ref{tab:appg-error} collects
the vocabulary. The book's own conclusion is a warning about reporting:
authors draw the cut differently, and a relative error measured at the training
resolution confounds the approximation term with the discretization term, so that if the
second is not small the printed number describes the quadrature, not the learned
operator. \index{error decomposition}

\item[\textbf{Filtered backprojection (FBP)}] The classical inversion
$f_{\text{FBP}} = \mathcal{R}^{*}\big(\mathcal{F}^{-1}(|\omega|\cdot \mathcal{F}p)\big)$
of \eqref{eq:ch36-fbp}, obtained by reading the Fourier slice theorem
\eqref{eq:ch36-slice} backwards: each projection is filtered in the detector variable
by the ramp $|\omega|$ and then smeared back along the rays that produced it by the
Radon adjoint $\mathcal{R}^{*}$. Exactness demands the polar sampling fill the Fourier
plane at the density the reconstruction grid requires --- complete angular coverage
and detector spacing below the Nyquist rate for the target pixel --- and the ramp is
a positive-order operator, gaining one derivative and amplifying high-frequency noise
exactly where the measured signal is weakest, so dose and resolution are traded
through the same factor. On sparse-view data the book is precise about what FBP is:
not an approximate inverse of the sampled operator $\mathcal{A} = \mathbf{M}
\mathcal{R}$, but the exact inverse of an operator whose input the scanner never
delivered, with the unwritten Fourier lines set to zero --- which is the origin of the
streak artifacts the learned reconstructors are scored against.
\index{filtered backprojection (FBP)|textbf}

\item[\textbf{Fourier multiplier}] The mode-indexed matrix
$R_{\phi}(k) \in \mathbb{C}^{d_{v} \times d_{v}}$ of \eqref{eq:ch05-multiplier} by
which the spectral branch scales the channel vector of the hidden field at each
retained frequency $k \in Z_{k_{\max}}$, and identically zero off it --- which is how
truncation enters the layer formula. The convolution theorem
\eqref{eq:ch05-convolution-theorem} is what makes the parameterization exact rather
than convenient: for a stationary kernel the exponentials are simultaneous
eigenfunctions and the eigenvalue at mode $k$ is the $k$-th coefficient of the
kernel, so learning $R_{\phi}$ and learning $\kappa(x-y)$ are the same object in two
coordinate systems, with $\kappa(x) = \sum_{k} R_{\phi}(k) e^{2\pi i k \cdot x}$
recovering the physical kernel. Channels mix inside a frequency and never across
frequencies; mode-to-mode transfer comes from the activation and the bypass, not the
multiplier. Conjugate symmetry \eqref{eq:ch05-reality} halves what must be stored,
and the book's implementation keeps the two low corners of the half-spectrum rather
than a centered box. \index{Fourier multiplier|textbf}

\item[\textbf{Fourier neural operator (FNO)}] The neural operator
of~\cite{li2021fourier} whose kernel layers are stationary convolutions on the torus,
written layerwise as \eqref{eq:fno_layer}: the activation of a sum of a local linear
map $W_{\ell} v$, the inverse transform of $R_{\phi,\ell}$ times the transform of $v$,
and a bias. Because the non-local part is diagonal in frequency, one layer costs
$\mathcal{O}(d_{v} N \log N)$ for the transforms plus a fixed contraction over the
retained modes, so the learned part of the computation stops growing with resolution;
that is the $\mathcal{O}(N \log N)$ entry of Table~\ref{tab:comparison}, whose per-layer
cost is the two-term count \eqref{eq:ch8-costs}.
Theorem~\ref{thm:fno_convergence} bounds its error on a bounded input set by
$C_{1} k_{\max}^{-\alpha} + C_{2}\mathcal{E}_{\text{NN}}(L, d_{v})$ for a target
gaining $\alpha$ derivatives, and \chapref{ch:part2-fourier-neural-operators-fno-from-scratch}
supplies the implementation while Appendix~\ref{app:b} supplies a runnable
one-dimensional case. The premise is bought, not free: stationarity is a claim about a
homogeneous medium with no distinguished location and periodic boundaries.
\index{Fourier neural operator (FNO)|textbf}

\item[\textbf{Fredholm integral kernel}] The trainable non-local ingredient of the
generic neural operator, $\big(\mathcal{K}_{\ell}(a; \phi_{\ell}) v_{\ell}\big)(x) =
\int_{\Omega} \kappa_{\ell}(x, y, a(x), a(y); \phi_{\ell})\, v_{\ell}(y)\, dy$ of
\eqref{eq:fredholm_kernel}. Its argument list is the point: the kernel sees the source
point, the target point, and the input field at both, which is what lets a layer
condition its own long-range interactions on the coefficient field and is exactly what
a convolution, whose kernel sees only $x - y$, cannot do. The motivation is the
representation Part~\ref{part:1} derives for elliptic problems, where the inverse of a
linear operator is an integral operator with a Green's kernel, suggesting that a
trainable kernel is the natural surrogate for a solution operator. The layer is also
the object to which quadrature consistency is owed: the integral is approximated by a
weighted nodal sum, and the sum must track the continuum integral as the mesh is
refined. \index{Fredholm integral operator|textbf}

\item[\textbf{Graph neural operator (GNO)}] The architecture of
\chapref{ch:part2-graph-neural-operators-gno-and-irregular-g} that keeps the generic
kernel layer \eqref{eq:no_layer} and restricts its \emph{support} rather than its
algebraic form: \eqref{eq:ch24-gno-kernel} integrates a kernel $\kappa_{\theta}(x_{i},
y)$ over the ball $\mathcal{B}_{r}(x_{i})$ and replaces the integral by a sum of
kernel values against the discrete measure $\sum_{j} w_{j}\delta_{x_{j}}$, so the mesh
enters only through the quadrature weights --- Voronoi dual areas, carried explicitly
because no constant spacing can absorb them. The implementation maps each edge
displacement $x_{i} - x_{j}$ through a shared multilayer perceptron to a channel-mixing
matrix and aggregates by scatter-add over the \texttt{edge\_index} neighbor list; a
radius-$r$ layer transmits nothing farther than $r$, which is why the geometry-informed
variant confines such layers to a point-to-grid encoder and a grid-to-point decoder and
runs the global part of the operator on a latent regular grid.
\chapref{ch:part1-comparative-theoretical-taxonomy-of-archit} records no convergence
rate and no discretization-invariance theorem for either construction, and the chapter
does not supply one. \index{Graph Neural Operator (GNO)}

\item[\textbf{Hilbert space}] A Banach space whose norm arises from an inner product
through $\|u\| = \sqrt{\langle u, u \rangle}$, equivalently a normed space satisfying
the parallelogram law; among the Lebesgue spaces $L^{p}(\Omega)$, only $p = 2$ qualifies,
and $H^{k}(\Omega) = W^{k,2}(\Omega)$ is Hilbert under \eqref{eq:ch02-hk-inner}. The
property the book uses is the isometric identification of the space with its dual:
Theorem~\ref{thm:ch02-riesz-rep} assigns each bounded linear functional a unique
representer of the same norm, so an error field and the error functional it induces are
two readings of one quantity, and orthogonal-projection reasoning --- the Fourier
partial sum \eqref{eq:ch05-projection} is exactly such a projection --- is available.
Banach spaces have no such identification, and that asymmetry is the reason negative-order
Sobolev spaces exist. \index{Hilbert space|textbf}

\item[\textbf{Hounsfield unit (HU)}] The linear attenuation scale on which computed
tomographic images are reported, and therefore the unit in which the reconstruction
error of \chapref{ch:part3-medical-imaging-and-tomographic-reconstruc} is stated.
The distinction matters to the book's argument about comparability: a root-mean-square
error in HU is a field error in the units of the measurement itself, so it is not on the
same footing as a normalized-image index, and a reported value wider than the whole
soft-tissue band reads as a reconstruction of raw sparse data rather than a tuned one.
\index{Hounsfield unit (HU)|textbf}

\item[\textbf{Input measure}] The Borel probability $\mu \in \mathcal{P}(\mathcal{A})$ on
the input space under which the population risk \eqref{eq:ch01-risk} takes its
expectation, and the reason a surrogate's guarantees are conditional. Training samples
are draws from it; the model inherits its support, so an error measured on smooth,
well-resolved coefficient fields licenses no claim about data from outside it, and the
book treats out-of-distribution failure as a property of the measure rather than of the
architecture. Two objects share the letter and are distinguished in the notation table:
the pointwise mean and standard deviation $\mu(x)$, $\sigma(x)$ of
\eqref{eq:ch21-normalize} are a Monte Carlo estimate of this measure's first two moments,
not the measure itself. On the theory side, Assumption~\ref{ass:ch09-compact} asks that
$\mu$ be supported on a set whose $\varepsilon$-entropy in the input norm is finite,
which is the statistical twin of the compactness hypothesis of
Theorem~\ref{thm:no_universality}. \index{input measure|textbf}

\item[\textbf{Lebesgue space $L^{p}$}] The space of measurable functions with
$\|u\|_{L^{p}} = (\int_{\Omega} |u(x)|^{p}\,dx)^{1/p}$ finite, $1 \le p < \infty$, with
the essential-supremum modification at $p = \infty$, as the notation table records;
$L^{p}(\Omega; \mathbb{R}^{d_{v}})$ is the vector-valued version for a field with $d_{v}$
components, and the flattening loss of \eqref{eq:rel_l2_loss} computes exactly that norm
with the channel count in place of $d_{v}$. Completeness is the Riesz--Fischer theorem, and only $p = 2$ is a Hilbert
space. In this book the exponent $p$ in $L^{p}$ and $W^{k,p}$ is always an integrability
exponent and never a parameter count or a latent width, as
\chapref{ch:part1-introduction-the-paradigm-shift-to-infinit} states and
Table~\ref{tab:notation-master} records. \index{Lebesgue space|textbf}

\item[\textbf{Lifting operator (lift)}] The first map $\mathcal{P} :
\mathbb{R}^{d_{a}} \to \mathbb{R}^{d_{v}}$ of the composition
\eqref{eq:no_composition}, a shallow multilayer perceptron applied pointwise to each
node's channels so that $v_{0}(x) = \mathcal{P}(a(x))$: a scalar permeability field
becomes a $d_{v}$-component latent field with no mixing of neighbors. Pointwise maps
can rescale and rotate channels but cannot carry information from $y$ to $x$, which is
why all non-locality is deferred to the kernel layers and why the division of labor
between the two is not cosmetic. In the implementation the lift is a $1 \times 1$
convolution sized for the augmented channel count \eqref{eq:ch21-augment}; in
Proposition~\ref{prop:ch06-zeroshot} it is the linear ingredient that preserves the band
limit, and so neither creates nor destroys the hypothesis the exact transfer statement
needs. \index{lifting operator|textbf}

\item[\textbf{Log-permeability measure}] The prior on the Darcy coefficient used by the
benchmark of \chapref{ch:part2-canonical-pde-benchmarks-and-reference-wor}: $\log K$ is
Gaussian with covariance operator $(-\Delta + \tau^{2} I)^{-d}$, $d = 2$,
\eqref{eq:ch25-darcy-measure}, a Mat\'ern-type covariance that is the standard prior for
subsurface coefficients. Read through its spectral decay, mode $k$ carries variance of
order $(1 + |k|^{2})^{-2}$, so sample paths of $\log K$ lie in $H^{s}$ for $s < 1$ and in
no better space; since continuity in two dimensions requires $s > d/2 = 1$, the
coefficient is almost surely unbounded, $K$ is log-normal rather than essentially
bounded, and Assumption~\ref{ass:ellipticity} fails on the support of the raw measure.
The book treats this as a specification to be repaired rather than a pathology to be
tolerated: draws are truncated into a band $K \in [K_{\min}, K_{\max}]$ before use, the
band is part of the benchmark, and $K_{\min}$ and $K_{\max}$ are to be reported as
carefully as the error. \index{log-normal permeability|textbf}

\item[\textbf{Mode truncation}] The restriction of the learned multiplier to the finite
box $Z_{k_{\max}} = \{k \in \mathbb{Z}^{d} : |k_{j}| \le k_{\max, j}\}$ of
\eqref{eq:ch05-modeset}, imposed by setting $R_{\phi}$ identically zero off the retained
set. A box rather than a ball, because it is a product of contiguous index ranges and a
batched transform can slice it without gathering scattered entries; the implementation
keeps two low corners of the half-spectrum, $2 m_{x} m_{y}$ modes against the
$(2m_{x}+1)(2m_{y}+1) \approx 4 m_{x} m_{y}$ of the full box, with the safeguard
$2 m_{x} < H$ and \texttt{modes2} held within $\lfloor W/2 \rfloor + 1$; exceeding the
second bound fails loudly as a shape error, exceeding the first lets the two slices
overlap and silently overwrite one another, which is the Nyquist-type hypothesis
$N_{\text{train}} > 2 k_{\max}$ of Proposition~\ref{prop:ch06-zeroshot} reappearing as a
sanity bound on a configuration file. The cost of truncation is the first term of
\eqref{eq:ch05-error-decomposition}, $C_{1} k_{\max}^{-\alpha}$: algebraic in the cutoff,
set by the smoothing order of the target and independent of depth, width, and grid.
Figure~\ref{fig:spectral_decay} computes the decay and the conclusion the book draws from
it --- once the algebraic curve falls below whatever the network's own error turns out to be,
a term this book measures for one architecture only, adding modes buys nothing, and a target with one derivative is not
rescued by any affordable budget, because regularity and not resolution is the binding
constraint. \index{spectral truncation|textbf}

\item[\textbf{Multi-query surrogate problem}] The name this book gives to the class of
workloads that motivate operator learning: settings in which the same governing
equations must be solved many times --- candidate geometries in shape optimization,
perturbed initial conditions in ensemble forecasting, sampled heterogeneity in reservoir
risk assessment --- so that ``the quantity of interest is not a solution but the map from
data to solution, and it must be evaluated at whatever resolution each query requires''.
The cost bookkeeping \eqref{eq:ch01-cost-bookkeeping} makes the term operational, giving
amortization only beyond the crossover $M^{\star}$, and Part~\ref{part:3} restates it with
the label cost made explicit: wherever the surrogate's per-query cost is below the
solver's, the break-even query count $Q^{\star}$ of \eqref{eq:ch31-ledger}, equivalently
\eqref{eq:ch38-crossover}, exceeds the number of labeled solves whenever the labels were
generated on a grid at least as fine as the one served. The gain the term names is the
amortization of a family of solves, not the removal of the resolution dependence. \index{multi-query surrogate problem|textbf}

\item[\textbf{Neural operator}] Any member of the parametric family
$\mathcal{G}_{\theta} : \mathcal{A} \to \mathcal{U}$ required to satisfy three conditions
stated in \chapref{ch:part1-introduction-the-paradigm-shift-to-infinit}: a parameter space
defined independently of any discretization, a continuum operator attached to each
parameter, and continuity in the norm of $\mathcal{U}$. Definition~\ref{def:generic_no}
supplies the concrete architecture, $\mathcal{G}_{\theta} = \mathcal{Q} \circ \mathcal{V}_{L}
\circ \cdots \circ \mathcal{V}_{1} \circ \mathcal{P}$ \eqref{eq:no_composition}, whose layers
$\mathcal{V}_{\ell}$ combine a local linear map, a bias, an activation, and a trainable
Fredholm kernel \eqref{eq:no_layer}. The book uses the phrase both for that architecture,
distinguishing it from DeepONet and the FNO, and for the class; Theorem~\ref{thm:no_universality}
establishes density for the class on compact subsets of $L^{2}$, requiring the target
operator to be continuous, and is silent on the parameter budget its proof implies.
\index{neural operator|textbf}

\item[\textbf{Operator learning}] The formulation in which the unknown is a map between
function spaces rather than a field. The contrast the book draws is with the
physics-informed neural network, where a network approximates one solution by minimizing a
residual at collocation points: there the unknown is a field and a new case is a new
optimization, whereas here the unknown is a map and a new input case is one forward
pass~\cite{kovachki2023neural, lu2021learning}. Its training objective is the expected
normalized error \eqref{eq:ch01-risk} under an input measure, and its targets are the
solution operators verified in \chapref{ch:part1-partial-differential-equations-as-solution}.
What the change of target forfeits is stated as plainly as what it buys: guarantees are
relative to a chosen space and a chosen compact subset of it, conserved quantities and
boundary conditions are not enforced by \eqref{eq:ch01-risk} unless the architecture or the
loss imposes them, and sample-complexity bounds under a general measure on function space
are not yet sharp. \index{operator learning|textbf}

\item[\textbf{Point evaluation functional}] The map $u \mapsto u(x)$, a linear functional
that is bounded on $H^{s}(\mathbb{T}^{d})$ exactly when $s > d/2$, by the coefficient bound
\eqref{eq:ch02-eval}, and unbounded on $L^{2}$ itself. That threshold decides which
architectures have a continuum interpretation at all: a branch network consuming nodal
values $u(x_{j})$ is reading continuous linear functionals only on spaces carrying more
than $d/2$ derivatives in the $L^{2}$ sense, and on $L^{2}$ its output has no well-defined
limit as the sampling is refined. The same condition legitimizes the restriction operator
$P_{N}$ of Definition~\ref{def:ch06-discretization-invariance}, and where inputs genuinely
are rough the book's repair is to read $P_{N}$ as a local average --- a cell mean, or a
moment against a narrow kernel --- which is the natural language of finite volume
discretization. \index{point evaluation functional}

\item[\textbf{Projection operator (project)}] The closing map $\mathcal{Q} :
\mathbb{R}^{d_{v}} \to \mathbb{R}^{d_{u}}$ of \eqref{eq:no_composition}, local and
pointwise, that returns hidden representations to physical solution channels through
$u(x) = \mathcal{Q}(v_{L}(x))$; like the lift it can rescale and rotate channels but
cannot move information between points, and in the implementation of
\eqref{eq:ch23-layer} it is two $1 \times 1$ convolutions with an activation between
them. The name is shared with two objects the book keeps distinct:
the $L^{2}$-orthogonal Fourier partial sum $P_{k_{\max}}$ of \eqref{eq:ch05-projection},
whose tail estimate \eqref{eq:ch02-tail} generates the truncation rate, and the
data-consistency update \eqref{eq:ch36-consistency} of
\chapref{ch:part3-medical-imaging-and-tomographic-reconstruc}, which acts on the
reconstructed image rather than on a channel index and, in the unrolled setting,
becomes a learned operator applied a fixed number of times.
\index{projection operator|textbf}

\item[\textbf{Pushforward regularization}] The training objective
$\mathcal{L}_{\text{pushforward}} = \sum_{k=1}^{K} \lambda_{k} \|\mathcal{G}_{\theta}^{k}(u_{0})
- u(k\Delta t)\|_{L^{2}}^{2}$ of \eqref{eq:ch26-pushforward}, which adds to the one-step
loss a weighted sum of the errors of the model's own iterates against solver snapshots,
with weights $\lambda_{k}$ and horizon $K$. Its mechanism is a change of the input
distribution and not a bound: no term in the objective contains a derivative of
$\mathcal{G}_{\theta}$, so it does not constrain the Lipschitz constant that governs
rollout growth, but it does place training inputs at the states the rollout visits, which
is where the finite differences forming the gradient are taken. The book states no
convergence theorem for it and notes that the ensemble statistics of
\eqref{eq:ch21-normalize} describe the solver's ensemble while a drifting trajectory is
measured against a reference the model is not drawing from. Applications use the same
device under other names: the two-step weather recipe is pushforward training ``in all
but name, at $K = 2$'', and the cardiac surrogate's loss \eqref{eq:ch07-data-loss} appears
optionally in this form. \index{pushforward regularization|textbf}

\item[\textbf{Quadrature consistency}] The requirement that a discrete nodal sum converge
to the continuum integral it stands in for, weights included; for the generic kernel layer
it is the substance of condition (iii) of Definition~\ref{def:ch06-discretization-invariance}
and of the argument behind \eqref{eq:limit_commute}, and for the periodic spectral layer it
is discharged by the trapezoidal rule, whose error is the aliasing identity
\eqref{eq:ch06-alias}. Its failure mode is named in \chapref{ch:part2-the-practitioner-s-paradigm-infinite-dimen}
in the strongest terms the book uses: on an irregular mesh there is no constant spacing to
absorb the weights, and if they are dropped ``the sum approximates nothing''. Consistency is
a statement about the discretization of a fixed continuum operator, and it is separate from
the invariance of the architecture family and from the accuracy of a trained model.
\index{quadrature consistency}

\item[\textbf{Quadrature weights}] The numbers $w_j$ against which a nodal sum is taken:
cell volumes or Voronoi dual areas, equal to $\Delta x^{d}$ on an equispaced grid, as the
notation table states. Two uses recur. In a kernel or graph layer they are what makes the
sum a discretization of an integral, and on a mesh they must be carried explicitly. In a
loss they cancel: for the relative $L^{2}$ metric \eqref{eq:rel_l2_loss} on equispaced
grids the same constant appears under both square roots, which is the precise and limited
sense in which that metric is invariant to grid density --- and the reason applying the
same unweighted code to a point-cloud tensor returns a quantity that is the $L^{2}$ norm
of nothing. \index{quadrature weights|textbf}

\item[\textbf{Relative $L^{2}$ error (metric)}] The quotient
$\mathcal{L}_{\text{rel-}L^{2}}(u, \hat{u}) = \|u - \hat{u}\|_{L^{2}(\Omega)} /
\|u\|_{L^{2}(\Omega)}$ of \eqref{eq:rel_l2_loss}, evaluated as a quadrature and the
implementation's root-then-average over the batch; it is the inner quotient of the
population risk \eqref{eq:ch01-risk} and the metric in which the benchmark percentages of
Part~\ref{part:2} and the measured error of Appendix~\ref{app:b} is reported, and being a
quotient of two norms it carries no physical units. Two qualifications are part of the
definition. Constant quadrature weights cancel, so the \emph{metric} is scale-free in $N$;
the \emph{value} still changes with $N$, because
$u$ and $\hat{u}$ are different functions sampled on different grids. And the guard added
to the denominator prevents division by zero but does nothing for conditioning when the
label norm is genuinely small, which the book handles by normalizing over a batch instead
of per sample. \index{relative $L^2$ error (metric)}

\item[\textbf{Rellich--Kondrachov theorem}] Theorem~\ref{thm:ch02-rellich}: for bounded
$\Omega \subset \mathbb{R}^{d}$ with $C^{1}$ boundary, the canonical injection
$\iota : H^{s}(\Omega) \hookrightarrow H^{t}(\Omega)$ \eqref{eq:ch02-injection} is a compact
linear operator whenever $s > t \ge 0$, and $H^{s}(\Omega) \hookrightarrow C^{k}(\overline{\Omega})$
continuously whenever $s > d/2 + k$. On the torus the proof is the mode-tail estimate
\eqref{eq:ch02-tail}: a bounded set in $H^{\sigma}$ lies within $\varepsilon$ of a bounded
set in a finite-dimensional space of band-limited functions, hence is totally bounded. The
theorem is the certificate the book trades a Sobolev regularity bound against, and the
reason the density theorems of Part~\ref{part:1} are stated on compacta: an operator such
as a viscous or elliptic solution map that gains $\alpha > 0$ derivatives sends bounded
sets to precompact ones. Its two limits are recorded where it is used --- it is qualitative, and
the smoothing order is a hypothesis about the target, not about the architecture.
\index{Rellich--Kondrachov theorem|textbf}

\item[\textbf{Resolution transfer}] Evaluation of a trained operator on a grid other than
its training grid, without retraining or interpolation of the output. The reading the book
insists on is narrower than the usual headline: the experiment recorded in
Appendix~\ref{app:b} reports relative $L^{2}$ error $0.0141$ at both the $64$- and the
$256$-point grid, a value quoted there and not plotted in
Figure~\ref{fig:super_resolution}, but its inputs carry eight Fourier modes and are already fully
resolved on the coarse grid, so the fine-grid evaluation resamples the \emph{same}
continuous functions. What is demonstrated is transfer of a learned operator, not recovery
of unresolved high-frequency detail, and not superiority over classical solvers on a linear
problem a single multiplier solves exactly. The governing distinction is stated once and
applied throughout: that a model \emph{can be evaluated} on any compatible grid is an
architectural property, while that it is \emph{accurate} there is an empirical claim which
must be re-verified per problem. \index{resolution transfer}

\item[\textbf{Restriction and extension operators}] The pair $P_{N}$ and $I_{N}$ linking a
continuum operator to its discrete realization: $P_{N} : a \mapsto (a(x_{1}), \dots,
a(x_{N}))$ records the nodal values a tensor stores, and $I_{N}$ reconstructs a function
from a discrete output, so that \eqref{eq:limit_commute} demands that down, across, and up
be asymptotically the same as straight across (Figure~\ref{fig:commutative}). The two array
forms of \eqref{eq:ch21-grid-layout} and \eqref{eq:ch21-node-layout} are the same $P_{N}$
written on a structured grid and on a point cloud. $P_{N}$ is legitimate only where point
evaluation is bounded, and the choice of $I_{N}$ is not innocent: for spectral layers the
book's canonical extension is the trigonometric interpolant, while re-interpolating an
output introduces an error the training loop never saw or bounded.
\index{restriction operator@restriction operator $P_N$|textbf}\index{extension operator@extension operator $I_N$|textbf}

\item[\textbf{\texttt{rfft2} (real two-dimensional FFT)}] The batched real-to-complex
transform \texttt{torch.fft.\allowbreak rfft2} applied over the two trailing spatial axes of a
channel-first grid tensor, the operation that makes the spectral branch one call rather
than a gather. Because a real field has conjugate symmetric coefficients
\eqref{eq:ch05-reality}, it returns only half the spectrum, of shape
$(B, C, N_{1}, \lfloor N_{2}/2 \rfloor + 1)$ in single-precision complex; the second line
of \eqref{eq:ch21-memory} prices that intermediate at eight bytes per entry, and
\texttt{torch.fft.\allowbreak irfft2} reconstructs the omitted half given the requested
output shape.
The spectral convolution slices its retained modes from this half spectrum in two
contiguous corners \eqref{eq:ch23-corners}, a union of blocks rather than a centered box,
so the retained set is defined by the transform's layout and not only by the wavenumbers
the model needs.
\index{rfft2 (real two-dimensional FFT)|textbf}

\item[\textbf{Root-mean-square error (RMSE)}] The square root of the mean squared nodal
difference, reported in two distinct units for two distinct purposes. In
\chapref{ch:part3-global-weather-forecasting-and-climate-dyn} it is quoted beside the
anomaly correlation coefficient precisely because it is the quantity an additive
displacement moves, and the published verification uses the area-weighted form matched to
\eqref{eq:ch33-accw}. In \chapref{ch:part3-medical-imaging-and-tomographic-reconstruc} it
is given in Hounsfield units as an image field error, \eqref{eq:ch36-metrics}. As a
\emph{training} loss the book rejects the unnormalized mean squared error it resembles: its
magnitude fluctuates with the array convention under refinement, and it lets one
large-magnitude sample dominate the objective, which is why the normalized quotient
\eqref{eq:rel_l2_loss} is used instead. \index{root-mean-square error (RMSE)|textbf}

\item[\textbf{Sensor (measurement location)}] One of the fixed points
$\{x_{j}\}_{j=1}^{m} \subset \Omega$ at which a DeepONet branch reads the input function,
assembled into $\mathbf{u}_{\mathrm{sensor}} = (u(x_{1}), \dots, u(x_{m})) \in
\mathbb{R}^{m}$. The set is ``the part of a DeepONet that the network never sees'': the
first linear layer consumes $m$ numbers in a fixed order, so the locations enter through
their count and their conditioning rather than through any learned geometry, and if two
ensemble members agree at the sensors while differing elsewhere the model is asked to
distinguish inputs its own argument list calls identical. The encoding term
$\mathcal{E}_{\text{sensor}}(m)$ measures exactly that ambiguity and decays like $m^{-s/d}$
for ensembles bounded in $H^{s}$, $s > d/2$, so halving it costs a factor $2^{d/s}$ in
readings --- regularity, not resolution, sets the price. Equispaced, space-filling, and
clustered node families differ along one axis: how density is distributed over $\Omega$,
and correspondingly which failure each guards against.
\index{sensors@sensor (measurement location)|textbf}

\item[\textbf{Sinogram}] The projection data set $p(\theta, s)$ of \eqref{eq:ch36-radon},
one axis the angle $\theta \in [0, \pi)$ and the other the detector position
$s \in [-R, R]$, so named because a point of the object traces a sinusoid through it.
Each entry is the line integral of the attenuation field along
$x_{1}\cos\theta + x_{2}\sin\theta = s$, and the idealization is where the first model
error enters: real beams are polychromatic, scatter adds a smooth positive bias, and photon
starvation where attenuation is high makes the noise neither Gaussian nor angle-independent.
Sparse-view CT samples the sinogram at $M \ll M_{\text{Nyq}}$ angles, leaving the sampled
operator an infinite-dimensional null space --- any field whose integral vanishes along
every measured ray is invisible to the data. \index{sinogram|textbf}

\item[\textbf{Sobolev embedding theorem}] Proposition~\ref{prop:ch02-sobolev-embeddings},
stated with its conditions exposed: on bounded $\Omega$ with $C^{1}$ boundary,
$W^{k,p} \hookrightarrow L^{q}$ continuously for $q$ up to the critical exponent
$p^{*} = dp/(d - kp)$ in the subcritical case $kp < d$, with compactness when $q < p^{*}$;
every finite $q$ at the borderline $kp = d$ but no $L^{\infty}$; and
$W^{k,p} \hookrightarrow C^{j,\alpha}$ when $j + \alpha \le k - d/p$ in the supercritical
case. On $\mathbb{T}^{d}$ the fractional statement is \eqref{eq:ch02-linfty-embed},
$H^{s} \hookrightarrow L^{\infty}$ if and only if $s > d/2$, and the condition is sharp ---
a bounded set in $H^{d/2}$ need not be bounded in $L^{\infty}$. The threshold is the book's
gate on everything pointwise: sensor readings, $L^{\infty}$ error claims, and the
Sobolev-regularized losses of \chapref{ch:part1-physics-informed-operators-and-sobolev-reg}.
\index{Sobolev embedding theorem|textbf}

\item[\textbf{Sobolev space}] The space $W^{k,p}(\Omega)$ of $L^{p}$ functions whose weak
derivatives $D^{\alpha} u$ through order $k$ are again in $L^{p}$, with the norm
\eqref{eq:ch02-wkp-norm}; Banach, and Hilbert when $p = 2$, where it is written $H^{k}$ and
carries the inner product \eqref{eq:ch02-hk-inner} that reads out all derivatives at once.
On the torus the scale extends to every real $s$ by the Fourier-coefficient condition
\eqref{eq:ch02-hs-torus}, so differentiability is polynomial decay of coefficients, negative
orders give the distributions dual to positive ones \eqref{eq:ch02-hsneg}, and the scale is
interpolating in the sense of \eqref{eq:ch02-interp}. Two members do most of the work in
this book: $H^{1}_{0}(\Omega)$, the closure of $C_{c}^{\infty}$, is the energy space of the
homogeneous Dirichlet problems and the codomain of the Darcy solution operator, and
$H^{s+\alpha}$ membership is the hypothesis that buys the truncation rate of
Theorem~\ref{thm:fno_convergence}. \index{Sobolev space|textbf}

\item[\textbf{Solution operator}] The map $\mathcal{G} : \mathcal{A} \to \mathcal{U}$ of
\eqref{eq:ch01-solution-operator} carrying data --- a coefficient field, an initial state,
a boundary forcing --- to a solution field, and the object this book learns rather than a
solution. The notation covers several distinct maps, and the differences matter: the
time-$t$ map $u_{0} \mapsto u(\cdot, t)$, the one-step map $u(t) \mapsto u(t + \Delta t)$
used for autoregressive rollouts, parameter-to-state maps in which finitely many physical
parameters enter as functions on $\Omega$, and the coefficient-to-state maps whose domain
and codomain differ and which are therefore typically smoothing. Before $\mathcal{G}$ exists
as an object of study, the three conditions of Hadamard well-posedness must be answered for
a specific \emph{pair} of norms --- existence, uniqueness, and continuous dependence ---
since the same evolution problem is Lipschitz between one choice of spaces and
discontinuous between another. \index{solution operator|textbf}

\item[\textbf{Spectral layer}] The FNO hidden layer \eqref{eq:fno_layer} as realized: two
parallel paths added before the activation, a local pointwise linear branch $W_{\ell}$ and a
global branch that transforms, multiplies the retained modes by $R_{\phi}$, and transforms
back, with a bias; Figure~\ref{fig:fno_spectral} and \eqref{eq:ch23-layer} give the
$(B, C, H, W)$ shapes at each stage. Three properties the book keeps separate: the global
branch is translation-equivariant and diagonal in frequency, so one layer may rotate the
channel vector within a mode but never moves energy between modes; the local bypass and the
activation are what break both diagonality and equivariance; and Parseval's identity makes
the operator norm of the spectral branch explicit as a supremum of the multiplier over the
retained modes. That per-layer constant is the ingredient of \eqref{eq:ch09-lipschitz},
whose product across layers is the global Lipschitz constant $\Lambda_{\theta}$ no training
objective in the book controls. Because the spectral and bypass paths meet
in an addition, a bypass demoted to half precision and cast back up recovers nothing.
\index{Fourier neural operator!layer definition|textbf}

\item[\textbf{Structural similarity index measure (SSIM)}] One of the three image-fidelity
numbers with which \chapref{ch:part3-medical-imaging-and-tomographic-reconstruc} frames
reconstruction quality, alongside the peak signal-to-noise ratio in decibels and the
root-mean-square error in Hounsfield units defined at \eqref{eq:ch36-metrics}; where those
two are functions of the pointwise difference, SSIM compares local luminance, contrast and
structure between the reconstruction and its reference over a sliding window. The book
adopts the name, states no formula of its own, and reproduces no operating point from the
earlier record, since the values there were measured on reconstructors the chapter
does not name and against one parameterization of the classical filter. Its substantive
argument about this family of metrics is a limitation: none of the three is a task-based
assessment of the diagnosis the image is meant to support, and SSIM additionally rewards
smoothness, so an over-smoothing network can score well while destroying the low-contrast
features the dose reduction was meant to preserve.
\index{structural similarity index measure (SSIM)|textbf}

\item[\textbf{Tensor layout}] The array convention by which a discretized field is stored:
the structured, channel-first grid layout $X \in \mathbb{R}^{B \times C \times N_{1} \times
\cdots \times N_{d}}$ of \eqref{eq:ch21-grid-layout}, and the point-cloud or mesh layout
$X_{\text{nodes}} \in \mathbb{R}^{B \times N_{\text{nodes}} \times C}$ with coordinates
carried separately, of \eqref{eq:ch21-node-layout}. The two differ along one axis: a weight
tensor in the first is indexed by lattice offsets and frequency bins, a weight tensor in the
second by the physical displacement $x_{i} - x_{j}$. The grid layout inherits translation
equivariance and the near-linear transform cost, and buys a kernel library; the mesh layout
inherits neither, aggregates over neighbor lists at a cost set by the edge count, and must
carry its quadrature weights explicitly. Only the batch axis carries no operator semantics,
which is why a finer grid is paid for out of it. \index{tensor layout}

\item[\textbf{Topological dual}] The space $X^{*} := \mathcal{L}(X, \mathbb{R})$ of
continuous linear functionals on a Banach space $X$, made Banach by the operator norm
\eqref{eq:ch02-dualnorm} even when $X$ itself is not complete, since a Cauchy sequence of
functionals converges pointwise. For a Hilbert space the dual is the space itself,
isometrically, by Theorem~\ref{thm:ch02-riesz-rep}; for $L^{p}$ the identification
$(L^{p}(\Omega))^{*} \cong L^{q}(\Omega)$ with $1/p + 1/q = 1$ holds only for
$1 \le p < \infty$ \eqref{eq:ch02-lpdual}, and the dual of $L^{\infty}$ is strictly larger
than $L^{1}$, being the space of bounded finitely additive measures absolutely continuous
with respect to Lebesgue measure. That failure is load-bearing: an $L^{\infty}$ bound on
permeability inputs yields no norm-convergent subsequence, and Banach--Alaoglu applied to
$L^{\infty} = (L^{1})^{*}$ buys only weak-$*$ compactness of the closed ball, which does
not control $\|a_{n} - a\|_{L^{\infty}}$. Reflexivity --- the canonical embedding of
$X$ into $X^{**}$ being onto, as for $L^{p}$ and $W^{k,p}$ with $1 < p < \infty$ --- is
the hypothesis that makes a bounded sequence yield a weakly convergent one.
\index{topological dual|textbf}

\item[\textbf{Torus}] $\mathbb{T}^{d} \cong [0,1)^{d}$ with opposite boundaries identified,
as the notation table defines it: the periodic domain on which the Fourier basis
\eqref{eq:ch02-fourier-pair} is exact, differentiation is multiplication by
$(2\pi i k)^{\alpha}$, and the convolution theorem diagonalizes every stationary operator at
once. It is the presupposition of the whole spectral construction, including the
$\mathcal{O}(N \log N)$ cost, and the three benchmarks accept it unevenly: Burgers' equation
and the Navier--Stokes vorticity form are posed periodically, which removes the boundary
condition from the specification and makes the pseudo-spectral reference solver cheap, while
Darcy flow carries homogeneous Dirichlet data on a square. Periodicity also replaces the
no-slip condition of a wall-bounded flow, and the book notes that this single choice limits
the benchmark's transfer to applications where boundary layers carry the drag.
\index{periodic domain (torus)|textbf}

\item[\textbf{Trunk network}] The DeepONet sub-network $y \mapsto \mathbf{t}(y) \in
\mathbb{R}^{p}$ that consumes output coordinates and supplies the second factor of
\eqref{eq:ch04-deeponet}; its $p$ outputs span at most a $p$-dimensional subspace, whatever
the branch learns, which is what the basis term $\mathcal{E}_{\text{basis}}(p)$ of
\eqref{eq:ch04-error-split} prices. Because the two halves are decoupled, one branch pass
serves every query point and the trunk is evaluated once per step when the query set is
shared across the batch \eqref{eq:ch22-batch-contraction}. The width $p$ is a rank
constraint fixed at construction; the residual output bias absorbs the ensemble mean of the
output field, and a stationary error floor as $p$ grows is a budget problem in the trunk
rather than an under-trained network. \index{DeepONet!trunk network|textbf}

\item[\textbf{Vorticity formulation}] The closed scalar form of the two-dimensional
incompressible Navier--Stokes equations, $\partial_{t}\omega + \mathbf{u}\cdot\nabla\omega =
\nu\Delta\omega + f$ with $\nabla \cdot \mathbf{u} = 0$ and
$\omega = \nabla \times \mathbf{u} = \partial_{1} u_{2} - \partial_{2} u_{1}$, posed on the
doubly periodic torus \eqref{eq:ch25-navierstokes}; the divergence-free velocity is
recovered from vorticity by the Biot--Savart law, so there is no pressure to eliminate. Two
specification points the book pins down. The forcing \eqref{eq:ch25-forcing} is Kolmogorov's
and, as written, is a scalar added to the vorticity equation --- the curl of a body force,
not a body force --- with $L^{2}(\mathbb{T}^{2})$ norm exactly $0.1$ and wavenumbers
$k_{2} = \pm 1$ only, so it injects energy at the largest scale. The viscosity range
$\nu = 10^{-3}$ to $10^{-5}$ is priced by the enstrophy balance \eqref{eq:ch03-enstrophy},
which at statistical steady state leaves $\nu\|\nabla\omega\|_{L^{2}}^{2} =
\langle f, \omega \rangle_{L^{2}}$ and therefore bounds the resolved gradient as viscosity
falls. \index{vorticity formulation|textbf}

\item[\textbf{Wall shear stress (WSS)}] The tangential traction
$\boldsymbol{\tau}_{w} = [\boldsymbol{\sigma}\hat{\mathbf{n}}]_{\text{tan}}$ exerted at a
boundary, equal for a Newtonian fluid to $\mu\, \partial_{n} u_{t}$ evaluated at the wall
\eqref{eq:ch32-wss}, with $\boldsymbol{\sigma}$ the viscous stress tensor, $\hat{\mathbf{n}}$
the outward unit normal, $u_{t}$ the tangential velocity, and $n$ the wall-normal direction.
It is the friction part of the force integral \eqref{eq:ch32-force-coeff}, whose other part
is pressure drag, and because it is a \emph{normal derivative} it is the quantity most
sensitive to the near-wall resolution that \eqref{eq:ch32-yplus} makes expensive. Two
consequences the book draws: a surrogate predicting pressure well and shear poorly can still
reproduce a bluff-body drag coefficient, so the two fields should be reported separately; and
separately predicted pressure and wall shear stress fields need not integrate consistently
with any divergence-free velocity. In the aerodynamic pipeline the surface WSS is among the
quantities the geometry-informed decoder reads back off the latent grid.
\index{wall shear stress (WSS)|textbf}

\item[\textbf{Weak derivative}] The function $v = D^{\alpha} u$ defined not pointwise but
against test functions, through $\int_{\Omega} u\, D^{\alpha}\phi =
(-1)^{|\alpha|} \int_{\Omega} v\, \phi$ for all $\phi \in C_{c}^{\infty}(\Omega)$
\eqref{eq:ch02-weakderiv}: integration by parts with the derivative transferred where it
costs nothing. Uniqueness almost everywhere follows from the density of the test space, and
the criterion the book holds onto is $u(x) = |x|$, whose weak derivative is $\operatorname{sign}(x)$:
kinks are admissible, jumps are not. That distinction propagates into the regularity
bookkeeping --- a one-dimensional jump has coefficients of size $|k|^{-1}$, so it lies in
$H^{s}(\mathbb{T})$ for $s < 1/2$ and in no better space, which is the functional-analytic
content of the Gibbs and shock-capturing difficulties recorded at the end of Part~\ref{part:1}.
\index{weak derivative|textbf}

\item[\textbf{Weak formulation}] The rewrite of a differential equation as an identity of
pairings --- find $u$ such that $B_{a}(u, v) = \langle f, v \rangle_{V^{*}, V}$ for every
test function $v$, as the Darcy weak form \eqref{eq:ch03-weakform} writes it, where
$B_{a}$ is the bilinear form of the coefficient $a$ and not the coefficient itself ---
with the pairing taken in a Gelfand triple $V \hookrightarrow H \cong H^{*} \hookrightarrow
V^{*}$ \eqref{eq:ch02-triple}, typically $H^{1}_{0}(\Omega) \subset L^{2}(\Omega) \subset
H^{-1}(\Omega)$, so that a source term that is a distribution rather than a function still
acts on the test space. The construction extends the inner product downward in regularity,
not upward, and its two analytic ingredients appear together: Poincar\'e's inequality
\eqref{eq:ch02-poincare} supplies the coercivity that the Lax--Milgram argument of
\chapref{ch:part1-partial-differential-equations-as-solution} invokes, and
\eqref{eq:ch02-energy} is the statement that testing the equation against the solution
itself is legitimate. \chapref{ch:part1-physics-informed-operators-and-sobolev-reg} turns the
same identity into a variational loss functional, which measures a residual in a weaker norm
than the $L^{2}$ error does. \index{weak formulation|textbf}

\item[\textbf{Zero-shot super-resolution}] The deployment literature's name for evaluating a
model trained on one grid on a finer grid without retraining, and the assertion the export
test of \chapref{ch:part2-production-deployment-and-high-throughput-} makes of a traced graph
by feeding it a fourfold finer tensor. The book uses the phrase when reporting others'
results --- evaluation of the aerodynamic surrogate on meshes sub-sampled by factors of two
to eight, where the geometry is fixed and only its sampling thins --- and reads it in every
instance through the caveat of
\chapref{ch:part1-discretization-invariance-and-continuum-li}: exact transfer holds for a
band-limited input whose spectrum lies inside the retained set, so what such experiments
demonstrate is resolution transfer on an already-resolved input rather than recovery of
frequencies absent from the training grid. The terms are not synonyms here, and the
difference is the point of the entry. \index{zero-shot super-resolution}

\end{description}

% NOTE: headwords and index keys were checked against every \index command in
% Parts I-III, and re-checked in the verification pass: no key carrying a textbf
% format here has a textbf anchor anywhere in the main text, so no glossary page
% duplicates a bold main-text entry. Six terms whose main-text top-level anchor is
% already bolded (resolution transfer, tensor layout, positional grid encoding,
% relative L^2 error, quadrature consistency, zero-shot super-resolution), three
% whose bold anchor sits on a subentry only (aliasing!quadrature error, point
% evaluation functional!boundedness threshold, error decomposition!statistical,
% truncation, and discretization), and the two terms opened and closed as ranged
% anchors in Ch. 1.8 (DeepONet, Graph Neural Operator (GNO)) therefore carry a plain
% index here. Keys reuse the main text's exact strings so that makeindex merges
% rather than splits them. The architecture-by-architecture inventory belongs to
% Appendix app:e, not to this chapter, and is not duplicated.
